\section{Exact coefficients and the composition interface}
\label{sec:composition-interface}

The proof begins with an exact coefficient, not with an estimate.  The finite Fourier identity below supplies its semantics.  The analytic argument is then separated into a local statement about the actual complex kernel and a weighted statement about everything outside the decoded well.  This separation is useful both for the present reciprocal system and for keeping clear which parts of the proof are classical and which are supplied by its arithmetic geometry.

\InputSectionBody{sections/02_fourier_selection}

\subsection{A source-twisted one-well composition principle}

Let $K\subset\mathbb R$ be compact and write
\[
 \Phi(u)=\frac{1}{\sqrt{2\pi}}\exp\!\left(-\frac{u^2}{2}\right).
\]
The next theorem packages only the final composition used later.  It does not assert that an observation or compression mechanism exists in an arbitrary system.

\begin{theorem}[Source-twisted one-well composition]
\label{thm:one-well-composition}
For each $X$, let $c_X(u)\geq0$ be an exact normalized coefficient, indexed by $u$ in a finite set $K_X\subset K$, and let $L_X>0$ and $\sigma_X\to0$.  Assume the following data.
\begin{enumerate}[label=\textup{(\roman*)}]
\item[\textup{(SEM)}] There is an exact finite-Abelian Fourier identity for $c_X(u)$, with the source twist $u$ retained in the actual complex summand.
\item[\textup{(OBS/COMP)}] An admissible provider partitions that exact sum into a decoded well and its complement.  The provider may be row elimination, anchor rigidity, decoder-image domination, product-fibre compression, or another mechanism proving the same two estimates below; none of these implementations is assumed universally.
\item[\textup{(LOCAL)}] For every fixed $C>0$, uniformly for $u\in K_X$,
\[
 L_X c_X^{\mathrm{well}}(u;C)
 =\frac{\Phi(u)}{\sigma_X}
  +O\!\left(\frac{\e^{-cC^2}}{\sigma_X}\right)
  +o_C(\sigma_X^{-1}),
\]
where the main term is obtained from the actual complex kernel on the decoded local well and $c>0$ is independent of $u$.
\item[\textup{(TAIL)}] After $C$ is fixed, the weighted contribution of every complementary lane satisfies
\[
 \sup_{u\in K_X}
 \left|L_X c_X^{\mathrm{tail}}(u;C)\right|
 =o_C(\sigma_X^{-1}).
\]
\item[\textup{(REAL)}] Positivity of the exact coefficient produces a witness in the finite quotient.  If an ambient equality is desired, an independent lifting or no-wrap criterion is supplied separately.
\end{enumerate}
Then, uniformly for $u\in K_X$,
\begin{equation}
\label{eq:one-well-composition-output}
 L_Xc_X(u)=\frac{\Phi(u)}{\sigma_X}+o(\sigma_X^{-1}).
\end{equation}
In particular, for all sufficiently large $X$ every source in the fixed compact window has positive coefficient and hence a quotient witness.  Whenever the independent lifting datum holds, the same witness realizes the intended ambient equality.
\end{theorem}

\begin{proof}
By SEM and the provider decomposition,
\[
 L_Xc_X(u)=L_Xc_X^{\mathrm{well}}(u;C)
             +L_Xc_X^{\mathrm{tail}}(u;C)
\]
exactly.  Insert LOCAL and TAIL, first fixing $C$ and then letting $X\to\infty$.  Since $K$ is compact, $\inf_{u\in K}\Phi(u)>0$.  Letting $C\to\infty$ gives \cref{eq:one-well-composition-output} uniformly and therefore positivity.  REAL converts positivity into a quotient witness; the final ambient equality uses only the separately stated lifting datum.  Thus $C$ is fixed first, then $X\to\infty$, and only afterwards is $C\to\infty$; no local phase is taken from the positive tail majorant.
\end{proof}

In the reciprocal-semiprime application, SEM is the sourced reciprocal coefficient identity; the provider is the one-anchor row-decoding and exact weighted fibre-compression mechanism; LOCAL is the actual-kernel Gaussian calculation; TAIL is the retained-factor six-sector estimate; and REAL is followed by the explicit no-wrap interval.  The remainder of this section develops the deterministic providers needed for that instantiation.

\InputSectionBody{sections/03_structural_tools}

The abstract principle therefore does not replace the arithmetic proof.  It records its interfaces, while the next section constructs the concrete reciprocal system on which those interfaces are verified.
